% ---------------------------------------------------------------- 00-front.tex
\begin{center}
{\LARGE\bfseries Same Pizza, Different Slices?\par}
\vspace{4pt}
{\large Putting the No-Tax Assumption of Modigliani--Miller Proposition I on Trial\par}
\vspace{10pt}
{\normalsize Lao Xiusheng, Xu Riqi, Wu Kaluk, Jiang Yumeng, Li Jiabao and Li Jingwen\par}
\vspace{2pt}
{\small\itshape Group D.E.\ Show\&Tell, MH6806 Principles of Finance and Risk Management (Topic 3)\par
Nanyang Technological University, October 2026\par}
\end{center}
\vspace{4pt}
\begin{center}\textbf{Abstract}\end{center}
\vspace{-6pt}
\begin{list}{}{\leftmargin=1cm\rightmargin=1cm\itemindent=0pt\listparindent=0pt\topsep=0pt\parsep=0pt}\item[]\small
Modigliani and Miller's Proposition I states that, in frictionless markets, the value of a firm is independent of its capital structure. We put a single assumption on trial, the absence of corporate income tax, and ask whether the interest tax shield turns debt into a free lunch. In a one-period, two-state economy, debt raises firm value by exactly the present value of the tax shield, but only up to the largest debt the weak state can service; beyond that point distress costs dominate. With a continuous cash-flow distribution the value function has an interior optimum, although the gain is small, under 1\% of unlevered value. What surprised us is that a higher tax rate makes each unit of debt more valuable but lowers the amount of debt that is safe, because the tax also comes out of the weak state's cash. We then ask whether this small gain reflects the mechanism itself or the simplifications of the setting. Letting lenders price default risk, limiting the usability of interest deductions and allowing debt to be rolled over all leave the structure of the trade-off intact; personal taxes can shrink the advantage or remove it, and the size of the gain depends on the horizon: with perpetual debt it rises to roughly 10\% of firm value. Applications to the 2017 US Tax Cuts and Jobs Act, the Singapore tax system and two Singapore-listed firms show that the model ranks debt capacity as the trade-off theory predicts, although its ``no debt'' verdict for the airline is driven by our no-cash-buffer default rule, and that observed leverage also reflects frictions outside the model. We conclude that the corporate tax makes financing matter without making debt free.\par
\smallskip
\textit{Keywords:} capital structure; Modigliani--Miller; interest tax shield; financial distress; trade-off theory\par
\textit{JEL classification:} G32, G33, H25
\end{list}
\vspace{6pt}
